%
%
%
This work presents a framework for studying temporal networks using zigzag persistence, a tool from the field of Topological Data Analysis (TDA). 
The resulting approach is general and applicable to a wide variety of time-varying graphs. For example, these graphs may correspond to a system modeled as a network with edges whose weights are functions of time, or they may represent a time series of a complex dynamical system. 
We use simplicial complexes to represent snapshots of the temporal networks that can then be analyzed using zigzag persistence. 
We show two applications of our method to dynamic networks: an analysis of commuting trends on multiple temporal scales, e.g., daily and weekly, in the Great Britain transportation network, and the detection of periodic/chaotic transitions due to intermittency in dynamical systems represented by temporal ordinal partition networks. 
Our findings show that the resulting zero- and one-dimensional zigzag persistence diagrams can detect changes in the networks' shapes that are missed by traditional connectivity and centrality graph statistics.

